% !TeX root = ../main.tex
\section{触发时间预测与时空预警（问题三）}\label{sec:warning}

\subsection{预报目标与时空定义}\label{sec:forecast-target}
% 定义源区触发时间、预警发布时间、预报窗口和下游到达时间。
\writingplaceholder

\subsection{基于动态风险的触发时间预报模型}\label{sec:time-forecast}
% 由第二问的动态触发风险进一步建立时间分布或离散时间风险模型。
\writingplaceholder

\subsection{触发时间的95\%预测区间构造}\label{sec:prediction-interval}
% 给出预测区间及计算方法，区分未来事件预测不确定性与模型参数不确定性。
\writingplaceholder

\subsection{预警阈值及预警提前量}\label{sec:alarm-lead-time}
% 根据预测风险制定报警规则，考虑观测、计算、发布等时间延迟。
\writingplaceholder

\subsection{区间覆盖率及预报可靠性}\label{sec:interval-coverage}
% 通过历史回测或独立样本评价预测区间覆盖情况；如果缺少足够历史事件，明确报告无法可靠估计实际覆盖率。
\writingplaceholder

\subsection{空间预警范围与传播距离估计}\label{sec:warning-extent}
% 结合 DEM 与沟道连通性估计可能影响的下游空间范围，并与第6章传播仿真结果进行衔接。
\writingplaceholder
